\section{Introduction}

Requirements engineering remains a central activity in software engineering because early misunderstandings about goals, constraints, stakeholders, and operating environments can propagate into architecture, implementation, testing, and deployment failures. Intelligent Requirements Engineering (IRE) investigates how knowledge-driven and multi-agent techniques can support requirements work with stronger traceability, validation, and reproducibility.

This paper studies a focused IRE module in the context of ChatREQ, a knowledge-driven multi-agent system for end-to-end requirements engineering. The module should be motivated by a concrete requirements engineering task, such as clarification, environment-aware generation, graph construction, conflict detection, human-in-the-loop adaptation, feasibility validation, architecture alignment, or executable task planning.

\subsection{Research Questions}

Replace the following placeholders with task-specific research questions:

\begin{itemize}
  \item RQ1: How effective is the proposed module compared with strong baselines?
  \item RQ2: Which knowledge sources, agent roles, or workflow components contribute most to performance?
  \item RQ3: How does the module affect human effort, traceability, downstream utility, or artifact quality?
\end{itemize}

\subsection{Contributions}

This paper should make explicit contributions such as:

\begin{itemize}
  \item A task formulation for an intelligent requirements engineering problem.
  \item A knowledge-driven method or multi-agent workflow.
  \item A dataset, benchmark task, or evaluation protocol.
  \item An empirical comparison against baselines and ablations.
  \item A reproducibility package with prompts, configurations, data, and scripts where permitted.
\end{itemize}

